\section{Methodology}

\subsection{PDE Solver}

I solve the one-dimensional advection-diffusion equation using the Finite Difference Method (FDM).

The advection term uses an upwind difference:

\[
advection = -v (u[i] - u[i-1]) / dx
\]

The diffusion term uses a central difference:

\[
diffusion = D (u[i+1] - 2u[i] + u[i-1]) / dx^2
\]

The solution is updated using:

\[
u\_new[i] = u[i] + dt (advection + diffusion)
\]

The initial condition is a Gaussian pulse:

\[
u = exp(-100 (x-0.3)^2)
\]

The grid uses 80 points, with:

\[
dt = 0.0005
\]

and:

\[
Nt = 150
\]

\subsection{Inverse Problem}

The goal is to find v and D that give a simulated solution close to the observed solution.

The loss function used in the code is:

\[
J(v,D) = ||u\_sim(v,D) - u\_obs||
\]

The optimization tries to make this value as small as possible.

\subsection{Genetic Algorithm}

I use a simple Genetic Algorithm (GA).

The initial population contains random values for v and D.

The algorithm selects the best solutions and creates new solutions using:

\begin{itemize}
    \item selection
    \item crossover
    \item mutation
\end{itemize}

The values are also kept inside the allowed parameter bounds.

\subsection{Particle Swarm Optimization}

I also use Particle Swarm Optimization (PSO).

Each particle has a position and a velocity. The position contains the values of v and D.

The particle updates its velocity using its own best position and the global best position.

The main parameters used are:

\[
w = 0.7,\quad c1 = 1.5,\quad c2 = 1.5
\]

PSO is used to search for parameter values with a small loss.